\section{Corrections to Prior Work and Recorded Dead Ends}
\label{sec:corrections}

Corrections and negative results delimit what the positive theorems actually
establish.  The first result below corrects a norm estimate in a publicly
posted QCPM complexity derivation.  The next two results isolate representation and
initialization costs on sparse Hamiltonian and central-path constructions.
We then record false shortcuts that were tested and rejected.  None of these
statements is a general verdict against Hamiltonian optimization,
preconditioning, or state reuse.

\subsection{Correction to the QCPM simulation norm}
\label{sec:corrections-qcpm}

We consider arXiv:2311.03977v2 (16 October 2024), the quantum central-path
method (QCPM) of Augustino, Leng, Nannicini, Terlaky, and Wu
\cite{augustino2024-quantum-central-path-algorithm-linear}.  All formulas in
this subsection refer to that version.  Let $d=m+n+2$ be the dimension of its
self-dual embedding, and use its notation
\begin{equation}
 s(z)=Mz+q,\qquad F(z)=z\odot s(z),\qquad
 f_\mu(z)=\frac12\norm{F(z)-\mu e}_2^2.                 \label{eq:corr-qcpm-potential}
\end{equation}
The central-path Hamiltonian is
$H(\mu)=-(h^2/2)\nabla^2+f_\mu$, and Algorithm~1 sets
% QCPM v2 has an internal normalization inconsistency between Algorithm 1's
% h(t) and equation (23); the definition of a here deliberately follows
% Algorithm 1, as do all constants below.
\begin{equation}
 h(t)=a\mu(t)^2,
 \qquad a=\frac{\gamma^2}{2R_1C_{d,\delta}},            \label{eq:corr-qcpm-h}
\end{equation}
where $C_{d,\delta}=\sqrt d/2+(3/4)\log(2/\delta)$ is the
concentration factor displayed in Algorithm~1.  We assume $d\geq1$,
$a,\eta>0$, and $0<\varepsilon<1$.  The spatial domain $\Omega_D$ is
fixed and compact and contains $e$; statements involving $z(\mu)$ also
require that point to belong to $\Omega_D$.  The invoked
real-space simulation theorem measures a potential $V$ on a fixed domain
$\Omega_D$ by
\begin{equation}
 \norm V_{\infty,1}=\int\sup_{z\in\Omega_D}|V(z,\tau)|\,d\tau.             \label{eq:corr-qcpm-simulator-norm}
\end{equation}
The QCPM proof truncates to a fixed box containing the central-path
neighborhood and replaces the supremum in
\eqref{eq:corr-qcpm-simulator-norm} by an assumed bound
$\sup_{\Omega_D} f_\mu\leq K\gamma^2\mu^2$.  That bound is false.

\begin{proposition}[Global-supremum counterexample]
\label{prop:corr-qcpm-supremum}
If $\Omega_D$ contains the initial center $e=z(1)$, then, for every
$0<\mu\leq1$,
\begin{equation}
 \sup_{z\in\Omega_D}f_\mu(z)
 \geq f_\mu(e)=\frac d2(1-\mu)^2.                       \label{eq:corr-qcpm-counterexample}
\end{equation}
Consequently, at $\mu=\varepsilon\leq1/2$, the asserted bound would require
$K\geq d/(8\gamma^2\varepsilon^2)$.
\end{proposition}

\begin{proof}
The self-dual embedding is initialized so that $s(e)=e$.  Thus $F(e)=e$,
and \eqref{eq:corr-qcpm-potential} gives
\[
 f_\mu(e)=\tfrac12\norm{(1-\mu)e}_2^2
          =\tfrac d2(1-\mu)^2.
\]
Taking the supremum proves \eqref{eq:corr-qcpm-counterexample}.  When
$\varepsilon\leq1/2$, its right-hand side is at least $d/8$, which gives the
claim about $K$.
\end{proof}

Ground-state concentration controls state-weighted mass, not the operator
supremum on a fixed domain.  A scalar phase shift does not restore the
claimed scaling.  If $z(\mu)\in\Omega_D$, then $f_\mu(z(\mu))=0$, and
for every real scalar $c$ the triangle inequality gives
\[
 \sup_{z\in\Omega_D}|f_\mu(z)-c|
 \geq\frac12|f_\mu(e)-f_\mu(z(\mu))|
 =\frac d4(1-\mu)^2.
\]

There is a second, independent issue in the time dilation.  Let
$\mu:[0,1]\to[\varepsilon,1]$ be continuous.  After division
by $h(t)\mu(t)$, the evolution equation is
\begin{equation}
 i\eta\,\partial_t\Psi
 =\left[-\frac{\theta(t)}2\nabla^2
       +\frac{f_{\mu(t)}}{\mu(t)h(t)}\right]\Psi,
 \qquad \theta(t)=\frac{h(t)}{\mu(t)}=a\mu(t).           \label{eq:corr-qcpm-scaled-evolution}
\end{equation}
The linear substitution used in the QCPM proof treats $\theta$ as constant,
although it varies with $t$.  The clock that normalizes the kinetic term is
instead
\begin{equation}
 \tau(t)=\frac1\eta\int_0^t\theta(u)\,du
        =\frac a\eta\int_0^t\mu(u)\,du.                 \label{eq:corr-qcpm-clock}
\end{equation}
Indeed, $dt/d\tau=\eta/\theta(t)$, so
\eqref{eq:corr-qcpm-scaled-evolution} becomes
\begin{equation}
 i\partial_\tau\widetilde\Psi
 =\left[-\frac12\nabla^2
        +\frac{f_{\mu(t(\tau))}}{h(t(\tau))^2}\right]
   \widetilde\Psi.                                      \label{eq:corr-qcpm-corrected-evolution}
\end{equation}
Combining \eqref{eq:corr-qcpm-h}--\eqref{eq:corr-qcpm-clock} yields the
exact corrected identity
\begin{equation}
 \boxed{
 \Lambda:=\norm V_{\infty,1}
 =\frac1{a\eta}\int_0^1
   \frac{\sup_{z\in\Omega_D}f_{\mu(t)}(z)}{\mu(t)^3}\,dt.}
                                                               \label{eq:corr-qcpm-exact-norm}
\end{equation}
This identity does not use a lower-bound argument.  The clock correction by
itself works in the authors' favor: it lowers the norm weight from the
constant-$\theta$ argument's $\mu^{-4}$ to $\mu^{-3}$.  The
$\varepsilon^{-3}$ blow-up below is driven entirely by the false global
supremum bound.

The schedule in Algorithm~1 is
\begin{equation}
 g(t)=c_e^{-1}\int_0^t
 \exp\!\left[-\frac1{\tau(1-\tau)}\right]d\tau,
 \qquad
 \mu(t)=\varepsilon+(1-\varepsilon)(1-g(t)),            \label{eq:corr-qcpm-schedule}
\end{equation}
where $c_e$ normalizes $g(1)=1$.
The bump integrand is extended by zero at both endpoints.  It is
continuous, nonnegative, and positive on $(0,1)$, so $c_e>0$.

\begin{lemma}[Width of the flat endpoint]
\label{lem:corr-qcpm-endpoint}
For $0<\varepsilon\leq\min\{1/4,c_e\}$, put
$u_\varepsilon=1/(2\log(1/\varepsilon))$.  Then
$\mu(t)\leq2\varepsilon$ for every
$t\in[1-u_\varepsilon,1]$.
\end{lemma}

\begin{proof}
For $0<u<1$, substitute $v=1-\tau$ in the tail of
\eqref{eq:corr-qcpm-schedule}.  Since $(1-v)v\leq v$,
\begin{equation}
 1-g(1-u)
 \leq\frac1{c_e}\int_0^u e^{-1/v}\,dv
 \leq\frac{u}{c_e}e^{-1/u}.                             \label{eq:corr-qcpm-tail}
\end{equation}
At $u=u_\varepsilon$, the exponential equals $\varepsilon^2$.  The last
quantity in \eqref{eq:corr-qcpm-tail} is at most $\varepsilon$:
$\varepsilon\leq1/4$ implies $0<u_\varepsilon\leq1/2$, and
$u_\varepsilon\varepsilon\leq c_e$ follows from $\varepsilon\leq c_e$.
For $t\geq1-u_\varepsilon$,
monotonicity then gives
$\mu(t)\leq\varepsilon+(1-\varepsilon)\varepsilon\leq2\varepsilon$.
\end{proof}

\begin{theorem}[Corrected QCPM simulator-norm scale]
\label{thm:corr-qcpm-norm}
For the choices in Algorithm~1 and
$0<\varepsilon\leq\min\{1/4,c_e\}$,
\begin{equation}
 \boxed{
 \Lambda\geq
 \frac{d}{128a\eta\,\varepsilon^3\log(1/\varepsilon)}.} \label{eq:corr-qcpm-lower}
\end{equation}
If $\Omega_D=[0,D]^d$, a valid global upper certificate is
\begin{equation}
 \Lambda\leq\frac{F_D}{a\eta\varepsilon^3},
 \qquad
 F_D=\frac d2\left[
 D\bigl(D\norm M_\infty+\norm q_\infty\bigr)+1
 \right]^2.                                             \label{eq:corr-qcpm-upper}
\end{equation}
Here $\norm M_\infty$ is the maximum absolute row sum.
\end{theorem}

\begin{proof}
On the interval supplied by \cref{lem:corr-qcpm-endpoint}, take
$\varepsilon\leq1/4$.  Proposition~\ref{prop:corr-qcpm-supremum} gives
$\sup_{\Omega_D} f_{\mu(t)}\geq d/8$, while
$\mu(t)^{-3}\geq(2\varepsilon)^{-3}$.  The interval has length
$1/(2\log(1/\varepsilon))$.  Substitution in
\eqref{eq:corr-qcpm-exact-norm} proves \eqref{eq:corr-qcpm-lower}.

For the upper bound, every $z\in[0,D]^d$ satisfies
\[
 |s_i(z)|\leq D\norm M_\infty+\norm q_\infty.
\]
Hence $|F_i(z)-\mu|$ is at most
$D(D\norm M_\infty+\norm q_\infty)+1$ for $0<\mu\leq1$, and
$f_\mu(z)\leq F_D$.  Since $\mu(t)\geq\varepsilon$, the exact identity
\eqref{eq:corr-qcpm-exact-norm} gives \eqref{eq:corr-qcpm-upper}.
\end{proof}

Since $a^{-1}=2R_1C_{d,\delta}/\gamma^2$,
\cref{eq:corr-qcpm-lower} gives the simulator-norm lower bound
\begin{equation}
 \Omega\!\left(
  \frac{R_1dC_{d,\delta}}
       {\gamma^2\eta\varepsilon^3\log(1/\varepsilon)}
 \right).                                                \label{eq:corr-qcpm-complexity-scale}
\end{equation}
By contrast, the bound claimed in v2 is
\[
 \norm V_{\infty,1}\leq\frac{K\gamma^2}{\theta\eta},
 \qquad
 O\!\left(\frac{K\gamma^2}{a\eta\varepsilon}\right)
 =O\!\left(\frac{R_1C_{d,\delta}}{\eta\varepsilon}\right)
 \quad(K=O(1)).
\]
Thus $\gamma^2$ cancels against $1/a$.  Conditional on all other simulation
hypotheses, \eqref{eq:corr-qcpm-upper} supplies the
norm-based upper certificate
\[
 O\!\left(\frac{F_DR_1C_{d,\delta}}
 {\gamma^2\eta\varepsilon^3}\right)
\]
before polylogarithmic factors and potential-evaluation cost.

The calculation has a schedule-independent form.

\begin{corollary}[Norm--speed tradeoff]
\label{cor:corr-qcpm-tradeoff}
Let $\mu:[0,1]\to[\varepsilon,1]$ be $L$-Lipschitz, with
$\mu(0)=1$, $\mu(1)=\varepsilon$, and $0<L<\infty$.
In particular, a differentiable schedule with $|\dot\mu|\leq L$
satisfies this hypothesis by the mean value theorem.  Retain $h=a\mu^2$.  For
$\varepsilon\leq1/4$,
\begin{equation}
 \boxed{\Lambda\geq\frac{d}{64a\eta L\varepsilon^2}.}   \label{eq:corr-qcpm-tradeoff}
\end{equation}
\end{corollary}

\begin{proof}
The endpoint values imply $L\geq1-\varepsilon$, hence
$\varepsilon/L\leq1$.  On the terminal interval
$[1-\varepsilon/L,1]$, the Lipschitz bound gives
$\mu(t)\leq\varepsilon+L(1-t)\leq2\varepsilon$.  Throughout that interval,
\cref{prop:corr-qcpm-supremum} gives
$\sup_{\Omega_D} f_\mu\geq d/8$, and
$\mu^{-3}\geq(2\varepsilon)^{-3}$.  Equation
\eqref{eq:corr-qcpm-exact-norm} proves the result.
\end{proof}
No monotonicity assumption is needed for this bound.

Taking $L$ large weakens \eqref{eq:corr-qcpm-tradeoff}, but moves $L$ into
the derivative, smoothness, and adiabatic estimates.  The corollary alone
does not optimize that tradeoff, and it does not exclude a theorem that
rigorously exploits localization.

As a numerical check, adaptive quadrature gives
$c_e=0.007029858406609657$.  For
\[
 J(\varepsilon)=\varepsilon^3\log(1/\varepsilon)
 \int_0^1(1-\mu(t))^2\mu(t)^{-3}\,dt,
\]
the computed values are
\begin{center}
\begin{tabular}{c@{\qquad}c}
\toprule
$\varepsilon$ & $J(\varepsilon)$ \\
\midrule
$10^{-2}$ & $0.77120646$ \\
$10^{-3}$ & $0.89055595$ \\
$10^{-4}$ & $0.95036732$ \\
\bottomrule
\end{tabular}
\end{center}
The sampled values are consistent with a limit of one; no limit is
deduced from this computation.  Retaining only the witness
$f_\mu(e)=d(1-\mu)^2/2$ gives
\[
 \Lambda\geq\Lambda_e
 :=\frac{d}{2a\eta}\int_0^1
 (1-\mu(t))^2\mu(t)^{-3}\,dt
 \geq\frac{d}{128a\eta\varepsilon^3\log(1/\varepsilon)}
 \quad\bigl(0<\varepsilon\leq\min\{1/4,c_e\}\bigr).
\]
The proof certifies $J(\varepsilon)\geq1/64$ in this range and makes no
claim that this constant is sharp.  The quadrature is a sanity check on the proved scale, not
an estimate of the full supremum or part of the proof.

\paragraph{Formal verification of the norm correction.}
The companion Lean development in \path{formal/QipmFormal/QCPM/}
checks the finite-dimensional potential and initialization identity,
the compact-domain supremum bound, robustness under scalar shifts,
the corrected clock and its inverse, and the integral change of variables.
It also checks the actual normalized bump schedule, the explicit
endpoint threshold, both lower-bound constants, and the fixed-box upper
certificate.  The headline bounds use the spatial supremum of the stated
potential; their continuity and integrability are proved rather than
assumed as extra properties of that supremum.  The initialization row
equations $Me+q=e$ and inclusion of the required points in the fixed domain
are explicit hypotheses.  The source correspondence is to Algorithm~1
of v2, whose normalization differs from its equation~(23).
The verification does not formalize Schr\"odinger evolution, the external
simulation or adiabatic theorems, or an unconditional query lower bound.

The norm and clock errors invalidate the stated complexity derivation, not
the QCPM idea itself, and \eqref{eq:corr-qcpm-complexity-scale} is not an
unconditional query lower bound for every implementation of the evolution.
Chakrabarti et al.\ analyze missing error controls in the pseudo-spectral
real-space simulation
\cite{chakrabarti2025-speedups-quantum-dynamics}.  Gribling et al.\ cite
that analysis and separately identify the absence of a suitable directly
applicable infinite-dimensional adiabatic theorem
\cite{gribling2025-self-concordant-schrodinger-operators-spectral}.  We do
not claim those broader criticisms as new.  Even after correcting the norm,
the following obligations remain:
\begin{itemize}[leftmargin=*]
 \item The schedule in \eqref{eq:corr-qcpm-schedule} is nonconstant and
 $C^\infty$-flat, but it is not analytic through either endpoint.  An
 analytic function with its zero endpoint Taylor series would be locally
 constant.  The quoted exponential adiabatic theorem assumes analytic
 continuation through the closed interval.
 \item The adiabatic theorem must be applied to the generator in
 \eqref{eq:corr-qcpm-scaled-evolution}; the derivative discussion analyzes a
 differently scaled Hamiltonian and does not establish uniform derivative
 constants for time-dependent $h$.
 \item The exact potential \eqref{eq:corr-qcpm-potential} is generally
 quartic.  The displayed Gaussian state and oscillator gap concern its
 quadratic approximation.  A dimension-, $\mu$-, and input-uniform transfer
 to the exact ground state and spectral gap is not established.
 \item The proposition and algorithm use different choices of $h$; the
 proposition's choice has no $\gamma$ or $\delta$ dependence and therefore
 cannot by itself provide arbitrary neighborhood opening and failure
 probability.
 \item The initial exact ground state is assumed rather than prepared.
 Knowing its minimizer $e$ does not prepare the ground state of a quartic
 Hamiltonian.
 \item The real-space theorem assumes a smooth periodic potential and
 periodic boundary conditions.  Restriction of
 \eqref{eq:corr-qcpm-potential} to a box neither makes it periodic nor
 controls the boundary and truncation errors.
 \item The box size must retain its solution-radius dependence unless an
 input-size promise bounds the central path.
\end{itemize}
A localized or clipped-potential repair would need leakage, spurious-state,
uniform-gap, boundary, and comparison estimates; concentration alone does
not supply them.

\subsection{Sparse metric, dense inverse metric}
\label{sec:corrections-densification}

The proposal of Gribling, Apers, Nieuwboer, and Walter replaces the Euclidean
Laplacian by the Laplace--Beltrami operator of the Hessian metric
\cite{gribling2025-self-concordant-schrodinger-operators-spectral}.  If
$G(x)=\nabla^2\phi(x)$, its principal part is
\begin{equation}
 L\psi=\frac1{\sqrt{\det G}}
 \sum_{i,j}\partial_i\!\left(
 (G^{-1})_{ij}\sqrt{\det G}\,\partial_j\psi
 \right).                                                \label{eq:corr-laplace-beltrami}
\end{equation}
Sparse constraints do not make these standard-coordinate coefficients
entrywise sparse.

\begin{theorem}[Path-simplex inverse metric]
\label{thm:corr-path-simplex}
Let
\[
 \mathcal D_n=\{x\in\R^n:0<x_1<\cdots<x_n<1\}
\]
and use the logarithmic barrier
\begin{equation}
 \phi(x)=-\log x_1-\sum_{i=1}^{n-1}\log(x_{i+1}-x_i)
          -\log(1-x_n).                                  \label{eq:corr-path-barrier}
\end{equation}
Every inequality contains at most two variables, and the
constraint-intersection graph is a path.  At the analytic center
$\bar x_i=i/(n+1)$, the Hessian $G=\nabla^2\phi(\bar x)$ is tridiagonal, but
\begin{equation}
 (G^{-1})_{ij}
 =\frac{\min(i,j)(n+1-\max(i,j))}{(n+1)^3}>0             \label{eq:corr-path-green}
\end{equation}
for every $i,j$.  For
$i,j\in[\lceil(n+1)/4\rceil,\lfloor3(n+1)/4\rfloor]$,
these entries are $\Theta(1/n)$, so $\Theta(n^2)$ coefficients are of
comparable order.
\end{theorem}

\begin{proof}
Write the $n+1$ slacks as
$s_0=x_1$, $s_i=x_{i+1}-x_i$ for $1\leq i<n$, and
$s_n=1-x_n$.  At $\bar x$, every slack is $1/(n+1)$.  If $B$ is the signed
path-incidence Jacobian of this affine slack map, then
\[
 G=B^T\diag(s^{-2})B=(n+1)^2T_n,
\]
where $T_n$ has diagonal entries two and adjacent off-diagonal entries
minus one.  The discrete Green identity is
\[
 (T_n^{-1})_{ij}
 =\frac{\min(i,j)(n+1-\max(i,j))}{n+1}.
\]
Scaling gives \eqref{eq:corr-path-green}.  In the middle half, both factors
in its numerator are bounded above by $n$ and below by a fixed positive
multiple of $n$, proving the final assertion.
\end{proof}

Thus a direct standard-coordinate, nondivergence-form finite-difference
stencil for the principal symbol in \eqref{eq:corr-laplace-beltrami} has
$\Omega(n^2)$ mixed-derivative directions at the analytic-center grid point,
although the input has only $O(n)$ nonzeros.  This is a densification warning,
not a simulation or query lower bound.  Indeed, at an arbitrary interior
point put $r_k=s_k^2$, $R=\sum_{k=0}^nr_k$, and let $i\leq j$.  Weighted path
inversion gives
\begin{equation}
 (G(x)^{-1})_{ij}
 =\frac{(\sum_{k=0}^{i-1}r_k)(\sum_{k=j}^{n}r_k)}{R}.     \label{eq:corr-path-semiseparable}
\end{equation}
Every strict off-diagonal block is rank one, and prefix sums give an $O(n)$
matrix--vector apply.  The witness therefore rules out only the inference
from sparse Hessian support to an entrywise sparse inverse-metric tensor.
It does not rule out implicit, transform-based, or coordinate-adapted
simulation.

Abe and Nagai parameterize their Riemannian Hamiltonian simulation by the
kinetic-matrix sparsity and separately assume sparse-access oracles for that
matrix; they note that constructing those black-box oracles is generally
nontrivial \cite{abe2026-quantum-riemannian-hamiltonian-descent}.
Theorem~\ref{thm:corr-path-simplex} shows only that input-level $O(n)$
sparsity is lost: the direct stencil has $\Theta(n^2)$ terms.  On a tensor
grid with $q$ points per coordinate it still has at most $2n^2+1$ neighbors
per row, sparse relative to the $q^n$-dimensional grid.  Thus the theorem
does not rule out sparse-Hamiltonian access or refute the separately assumed
Abe--Nagai oracles.  Gribling et al.\ likewise leave the simulation and
qubit-realization questions open.

\subsection{A hidden-sign start-oracle separation}
\label{sec:corrections-hidden-sign}

We reuse, without reconstructing it, the path-sparse hidden-sign LP
\eqref{eq:structural-locality-lp} from
\cref{sec:structural-geometric-locality}.  Recall that
$p_0=1$, $p_i=\prod_{j=1}^i\varsigma_j$, and
$q_i=u_i+v_i$.

\begin{theorem}[Value, output, and central-start separation]
\label{thm:corr-hidden-sign}
For the family \eqref{eq:structural-locality-lp}:
\begin{enumerate}[leftmargin=*]
 \item It is strictly primal--dual feasible and has the unique optimum
 \begin{equation}
  \operatorname{OPT}=N+1+\alpha p_N.                    \label{eq:corr-hidden-opt}
 \end{equation}
 Estimating its shifted value $\operatorname{OPT}-(N+1)$ to additive error
 less than $\alpha$ requires $\Omega(N)$ bounded-error raw sparse queries.
 \item Every classical nonnegative output $\widehat x$ satisfying
 $\norm{A_\varsigma\widehat x-b}_1<1/2$ determines $p_N$ and therefore also
 requires $\Omega(N)$ queries.
 \item The primal central path has the public support-two orthonormal null
 basis stated in \cref{sec:structural-geometric-locality}, and its reduced
 log-barrier Hessian has condition number one for every $\mu>0$.
 \item At $\mu=3/4$, one lookup at the final cell of a reusable random-access
 value oracle for the exact primal central point reveals $p_N$.  If setup and
 that lookup use respectively $Q_{\rm setup}$ and $Q_{\rm lookup}$ raw
 queries, then
 $Q_{\rm setup}+Q_{\rm lookup}=\Omega(N)$.  In particular, query-free
 invocations after preprocessing force $Q_{\rm setup}=\Omega(N)$.
\end{enumerate}
\end{theorem}

\begin{proof}
The feasibility argument and central-path calculations are given throughout
\cref{sec:structural-geometric-locality}, with the displayed formulas in
\eqref{eq:structural-locality-barrier}--
\eqref{eq:structural-locality-condition-one}.  In particular, feasibility
fixes $d_i=p_i$ and leaves $q_i\geq1$, so the objective is
$\sum_iq_i+\alpha p_N$ and has the unique minimum
\eqref{eq:corr-hidden-opt}.  The shifted value reveals parity.

For an approximate output define
$r_0=\widehat d_0-1$ and
$r_i=\widehat d_i-\varsigma_i\widehat d_{i-1}$.  Telescoping gives
\begin{equation}
 \widehat d_N-p_N
 =\sum_{j=0}^N\left(\prod_{k=j+1}^N\varsigma_k\right)r_j,               \label{eq:corr-hidden-telescope}
\end{equation}
so $|\widehat d_N-p_N|<1/2$ under the asserted residual bound.  The sign of
$\widehat d_N=\widehat u_N-\widehat v_N$ is therefore $p_N$.

Finally, \eqref{eq:structural-locality-center} gives $q=2$ at $\mu=3/4$.
The final pair is $(3/2,1/2)$ or its swap according to $p_N$, so one
coordinate-oracle lookup determines parity even with coordinate error below
$1/2$.  Since quantum parity has query complexity $\Theta(N)$, each claimed
raw-query lower bound, including the combined setup--lookup bound, follows
from the local sign-oracle reduction used throughout \cref{sec:lp-hardness}.
\end{proof}

This theorem charges an original-form feasible start.  It does not apply to
infeasible-start or homogeneous self-dual methods, which can supply a
different universal start; nor does it say that the scalar reduced Newton
solve itself is hard.

\subsection{Dead ends and false shortcuts}
\label{sec:corrections-dead-ends}

\paragraph{First inverse versus filtering.}
A small value of
$\norm H\norm{H^{-1}f}/\norm f$ does not imply a fast DLS filtering solve.
Filtering depends on the second inverse through
$\norm H\norm{H^{-2}f}/\norm{H^{-1}f}$, and active--inactive coupling can
restore the full condition-number exponent.  The exact identities and sparse
sharpness pair are in
\cref{sec:beyond-kappa-coupling,sec:beyond-kappa-truncation}.

\paragraph{Sparse support versus inverse access.}
Sparse Hessian or KKT support does not imply a sparse inverse, a cheap block
encoding after elimination, or low circuit depth on bounded-range hardware.
The factor-normalization and inverse-compilation frontiers in
\cref{sec:preconditioning-whitening} prove the access cost; the bounded-range
depth claim requires the bounded-occupancy geometric input placement of
\cref{sec:structural-geometric-locality}.  The dense metric in
\cref{sec:corrections-densification} gives a separate coordinate-level
warning.  None of these statements says that every structured inverse is
expensive.

\paragraph{Condition one versus the surrounding maps.}
A condition-one reduced Newton matrix does not make the affine feasible
offset, coordinate gauge, transformed right-hand side, recovery map, or
requested output free.  The LP separations in
\cref{sec:lp-affine-separation,sec:lp-prefix-rigidified} place parity in
the affine offset while leaving the reduced geometry public.  The full-KKT
and preconditioner dichotomies in
\cref{sec:lp-full-kkt,sec:preconditioning-parity} show how input dependence
can relocate among these objects.

\paragraph{Exact Schur preconditioning.}
An exact Schur transformation can make the abstract product constantly
conditioned while relocating input dependence into construction, transformed
RHS loading, or recovery.  The complete ledger is proved in
\cref{sec:lp-full-kkt,sec:preconditioning-whitening}.  Supplying the
preconditioned product itself as a unit-cost oracle is a strictly stronger
input model, not a consequence of sparse access to the original system.

\paragraph{State output versus coordinate access.}
A normalized solution state is not a reusable coordinate oracle.  Its norm
and global phase are absent, while coherent diagonal construction also needs
controlled and adjoint access and pays a normalization factor.  The tight
state-to-diagonal product appears in
\cref{sec:preconditioning-state-interface}; the implicit-coordinate escape
route in \cref{sec:positive-compressed-dual} states the stronger interface it
actually constructs.

\paragraph{One-step bounds versus a fixed trajectory.}
Reapplying a one-step lower bound at every point of one fixed LP trajectory
is invalid without freshness, a memory restriction, an adaptivity/causality
restriction, or mandatory explicit output.  The fresh-oracle direct sum is
\cref{thm:prec-online-direct-sum}; the surrounding discussion shows how
caching a static input can erase later raw-query cost.

\paragraph{Residual bounds versus projective variation.}
Ordinary short-step neighborhood and relative OSS residual bounds do not
imply finite projective variation.  Section~\ref{sec:positive-lazy} gives a
sparse cycling counterexample and the additional hypotheses needed for
reuse.  The refined audit in \cref{sec:positive-face-repair} also corrects
an earlier version's proposed sharpness family: its convergent leakage and
transverse-forcing ledgers still allow linearly growing realized variation,
and a separate oscillating-displacement family shows that adding the
$O(\mu_t)$ tail bound on face leakage does not help either; finite realized
variation is therefore stated as an independent hypothesis rather than
derived from the ledgers.

\paragraph{Expanderization.}
Replacing a path by an expander does not evade the bounded-degree
local-linear mixture or cut/resistance obstruction.  Those no-go results are
\cref{sec:frontier-mixture,sec:frontier-electrical}.  In the favorable
well-spread regime, \cref{sec:positive-checkpoint-span} gives the matching
classical sampling shortcut; making rare columns influential instead moves
cost into normalization or one-time search.

\paragraph{Scalar gap versus density visibility.}
A constant scalar objective gap can coexist with a trace-normalized central
density that is only $O(1/G)$ from a public state.  Thus an output lower bound
must state its accuracy and normalization contract.  The zero-query boundary
and the mass--visibility inequalities are proved in
\cref{sec:condensation-visibility,sec:condensation-query-lower}; the dilution
example is \cref{rem:wta-dilution}.

\paragraph{Pairwise distance versus a common decoder.}
Pairwise distinguishability of states indexed by an unknown winner set does
not provide one public measurement that decodes every set.  The failed
collision shortcut loses a factor $1/k$.  Section~\ref{sec:state-conversion-model}
repairs this only by measuring a proposed label and verifying that label with
fresh input queries; the verifier, not pairwise trace distance alone, supplies
the common decision procedure.

\subsection{Disposition of remaining algorithmic candidates}
\label{sec:corrections-dispositions}

To state the factorized rectangular observation precisely, let $B_\mu$ have
full row rank, let $f=B_\mu h$ lie in $\operatorname{range}(B_\mu)$, and let
$\alpha_{B_\mu}\geq\norm{B_\mu}$ be its block-encoding normalization.  Define
\begin{equation}
 \rho_{B_\mu}(f)=\alpha_{B_\mu}
 \frac{\norm{(B_\mu B_\mu^T)^{-1}f}_2}
 {\sqrt{f^T(B_\mu B_\mu^T)^{-1}f}}.                    \label{eq:corr-rectangular-rho}
\end{equation}
This pseudoinverse parameter is distinct from the nonsingular filtering
parameter in \eqref{eq:dls-general-filtering-parameter}.  In the mixed
strictly-complementary exact-centering asymptotic, assume
$B_\mu B_\mu^T=\mu^{-1}C+\mu D+o(\mu)$,
$f=b\in U:=\operatorname{range}C$, and the generic active--inactive coupling
of \cref{sec:beyond-kappa-coupling}, so that
\[
 (B_\mu B_\mu^T)^{-1}b=\mu(p,-g)+o(\mu),
 \qquad b^T(B_\mu B_\mu^T)^{-1}b=\Theta(\mu).
\]
With the essential normalization
$\alpha_{B_\mu}=\Theta(\mu^{-1/2})$, substitution in
\eqref{eq:corr-rectangular-rho} gives the proved algebraic observation
$\rho_{B_\mu}(b)=\Theta(1)$.

Here $D_\mu=(XS^{-1})^{1/2}$, $B_\mu=AD_\mu$, and
$h=(XS)^{-1/2}r_c$ is the scaled complementarity right-hand side; thus
$f=B_\mu h$ is the row-space right-hand side denoted ``$Bh$''.  The two
scaled direction components are
\[
 u=\Pi_{\ker B_\mu}h,
 \qquad
 v=B_\mu^+f=B_\mu^+(B_\mu h),
\]
with $u+v=h$.  They recover the physical directions through
$\Delta x=D_\mu u$ and $\Delta s=D_\mu^{-1}v$.  The
constant value of \eqref{eq:corr-rectangular-rho} is not an end-to-end QIPM
speedup.  Six obligations remain: a DLS-quality rank-deficient rectangular
filter theorem with its cutoff and error law; preparation of $f=B_\mu h$;
coherent construction, norm control, and cancellation-safe physical recovery
of $u$ and $v$; inexact-neighborhood propagation; maintenance of the changing
diagonal oracles; and a complete output ledger with a strong classical
comparison.  Li's instance-dependent rectangular minimum-norm algorithm
already uses closely related second-inverse geometry
\cite{li2025-new-condition-number-application-multivariate-polynomial}, so
the primitive also substantially collides with prior work.

A further candidate, the augmented spectral Newton sketch, instead samples the
augmented matrix \([D_{\rm N}^{1/2}A^T,v_{\rm N}]\), where \(D_{\rm N}\) and
\(v_{\rm N}\) denote the sketch's Newton scaling and regression target.
Under its stated row-access assumptions,
a \((1\pm\varepsilon)\) embedding of the augmented Gram matrix gives a valid
least-squares embedding lemma at
\(\Otilde(\sqrt{nm}/\varepsilon)\) row-query scale; the infeasible-step
analysis, implicit update mechanism, and end-to-end comparison remain open,
and the direction overlaps the factorized rectangular candidate above.

State-to-diagonal lifting is a useful but known primitive, not a new result:
Rattew and Rebentrost give the amplitude-to-diagonal block encoding, with the
phase-aligned coherent access assumptions recorded in
\cref{sec:preconditioning-state-interface}
\cite{rattew2023-non-linear-transformations-of-quantum}.  Likewise, the
projective-plus-implicit expander candidate is closed by the dichotomy in
\cref{sec:positive-checkpoint-span}: bounded contributions admit a comparable
classical sketch, while rare influential columns incur normalization or
one-time search cost.
